\section{表格和图形}

\subsection{表格}

使用 tabular 环境, 如下表格: 表~\ref{tab:foo}. 通过 \verb|autoref| 引用表格: \autoref{tab:foo}.

这里可使用命令 \verb|\tabcolsep| 和 \verb|\arraystretch| 分别控制列间距和行间距.

\begin{table}[htp!]
\centering
\setlength{\tabcolsep}{12pt}  % 6pt standard
\renewcommand{\arraystretch}{1.2}
\caption{学术活动安排样例}
\label{tab:foo}
\begin{tabular}{|c|c|c|c|}
\hline
\textbf{日期}  & \textbf{地点} & \textbf{活动名称} & \textbf{备注} \\ \hline
2024年8月1日      & 上海       & 学术研讨会      & 主题：人工智能 \\ \hline
2024年8月15日    & 北京       & 学术交流会      & 重点：数学建模 \\ \hline
2024年9月1日      & 深圳       & 研究研讨会      & 主题：数据科学 \\ \hline
2024年10月15日  & 广州       & 创新论坛         & 重点：科技创新 \\ \hline
\end{tabular}
\end{table}


\clearpage

三线表
\begin{table}[htp!]
\newcolumntype{L}{X}
\newcolumntype{C}{>{\centering \arraybackslash}X}
\newcolumntype{R}{>{\raggedright \arraybackslash}X}
\centering
\caption{某校学生升高体重样本}
\label{tab:heightweight}
\begin{tabularx}{0.9\textwidth}{CCCC}
\toprule
序号 & 年龄 & 身高 & 体重 \\
\midrule
001 & 15 & 156 & 42 \\
002 & 16 & 158 & 45 \\
003 & 14 & 162 & 48 \\
004 & 15 & 163 & 50 \\
\cmidrule{2-4}
平均 & 15 & 159.75 & 46.25 \\
\bottomrule
\end{tabularx}
\end{table}

% 某行业产量与生产费用的数据
% \begin{table}[htp!]
% \centering
% \caption{某行业产量与生产费用的数据}%\label{}
% \newcolumntype{Y}{>{\centering\arraybackslash}X}
% \newcolumntype{Z}{!{\vline}@{\color{white}\vrule width \doublerulesep}!{\vrule}}%自定义列格式(双线)
% \begin{tabularx}{0.94\textwidth}{c|c|YZc|c|Y}
% \Xhline{0.9pt}
% 企业编号&  产量(台)&生产费用(万元)&企业编号&产量(台)&生产费用(万元)\\\Xcline{1-3}{0.6pt}\Xcline{4-6}{0.6pt}
% 1 & 40 & 130 & 7&  84&  165\\
% 2 & 42 & 150 & 8&  100&  170\\
% 3 & 50 & 155 & 9&  116&  167\\
% 4 & 55 & 140 & 10&  125&  180\\
% 5 & 65 & 150 & 11&  130&  175\\
% 6 & 78 & 154 & 12&  140&  185\\\Xhline{0.72pt}
% \end{tabularx}
% \end{table}

研究生数学建模2019年F题结果示例
\begin{table}[htp!]
\centering
\caption{问题1结果1 (左) 与 问题2结果 (右)}
\begin{minipage}[h]{0.48\linewidth}
\renewcommand\arraystretch{1.2} %定义表格高度
\newcolumntype{Y}{>{\centering\arraybackslash}X}
\begin{tabularx}{0.9\textwidth}{|Y|Y|Y|}
  \hline
  数据集1  &  数据集1 & 数据集2  \\
 \hline
  A问题1   & A问题1 &  A问题1     \\
  \hline
  503   &  503     & 163     \\
  294   &  200     & 114      \\
  91    &  80      & 8     \\
  607   &  237     & 309      \\
  540   &  170     & 305    \\
  250   &  278     & 123    \\
  340   &  369     & 45      \\
  277   &  214     & 160    \\
  B     &  397     & 92    \\
        &  B       & 93    \\
        &          & 61        \\
        &          & 292       \\
        &          & B         \\
104861  & 103518   & 109342 \\
\hline
\end{tabularx}
\end{minipage}
\begin{minipage}[h]{0.48\linewidth}
\renewcommand\arraystretch{1.2} %定义表格高度
\newcolumntype{Y}{>{\centering\arraybackslash}X}
\begin{tabularx}{0.9\textwidth}{|Y|Y|Y|}
  \hline
  数据集1  &  数据集1 & 数据集2  \\
  \hline
  A问题2  &  A问题2  & A问题2    \\
  \hline
   503    & 503     & 163  \\
  294    & 200     & 114   \\
   91     & 80     & 8   \\
   607    & 237     & 309   \\
   540    & 170    & 305  \\
   250    & 278    & 123  \\
   340    & 369     & 45   \\
   277    & 214    & 160  \\
  B      & 397   & 92 \\
        &  B     &  93      \\
         &         &  61      \\
        &         &   292     \\
        &         &   B     \\
 104917  &  103563  &109427 \\
\hline
\end{tabularx}
\end{minipage}
\end{table}

\clearpage
\begin{table}[htp!]
%\small
\centering
\renewcommand\arraystretch{1.2} %定义表格高度
\newcolumntype{Y}{>{\centering\arraybackslash}X}
\caption{问题3结果}
\begin{tabularx}{0.9\textwidth}{|Y|Y|Y|Y|Y|Y|}
  \hline %定义表格宽度
  数据集1  &  数据集1  &  \multicolumn{2}{c|}{数据集2 (无问题点)}  & \multicolumn{2}{c|}{数据集2 (有问题点)} \\
  \hline
  A问题3   & A问题3   &  \multicolumn{2}{c|}{A问题3}    &  \multicolumn{2}{c|}{A问题3}      \\
  \hline
  503      & 503     & 169      & 73      & 169     & 73   \\
  \hline
  69       & 69      & 322      & 249     & 322     & 249   \\
  \hline
  506      & 506      & 270     & 274     & 270     & 274   \\
  \hline
  371      & 371      & 89      & 12      & 89      & 12   \\
  \hline
  183      & 183      & 236     & 216     & 236     & 216  \\
  \hline
  194      & 194      & 132     & 16      & 132     & 16   \\
  \hline
  450      & 450      & 53      & 282     & 53      & 282   \\
  \hline
  286      & 113      & 112     & 84      & 112     & 141  \\
  \hline
  485      & 485      &  268    & 287     & 268     & 291 \\
  \hline
 \red{B (9D)}~~     & 248      & 250     & 99      & 250     &161 \\
 \hline
   & \red{B (10D)}   & 243     & \red{B (21D)}   & 243    & \red{B (21D)} \\
   \hline
    &          &         &         &         &       \\
    \hline
  104861m  & 103518m  &         & 168924m  &         &161650m \\
\hline
\end{tabularx}
\end{table}


\subsection{图形}

图形并列
\begin{figure}[htp!]
\begin{minipage}[t]{0.48\linewidth}
\centering
\includegraphics[width=0.9\textwidth]{image1.eps}
\caption{fig1}
\label{fig:side:a}
\end{minipage}%
\begin{minipage}[t]{0.48\linewidth}
\centering
\includegraphics[width=0.9\textwidth]{image2.eps}  % 2.2in
\caption{fig2}
\label{fig:side:b}
\end{minipage}
\end{figure}

这是一个算法流程图
\begin{figure}[htp!]
\centering
\includegraphics[width=.55\textwidth]{fig.png}
\caption{算法流程图}
\end{figure}

多图并排
\begin{figure}[!htp]
	\centering
	\subfloat[Arabic numerals]{\includegraphics[width=0.4\textwidth]{image1.eps}}\qquad
	\subfloat[Arabic numerals]{\includegraphics[width=0.4\textwidth]{image2.eps}} \\
	\subfloat[Arabic numerals]{\includegraphics[width=0.4\textwidth]{image3.eps}}\qquad
	\subfloat[Arabic numerals]{\includegraphics[width=0.4\textwidth]{image4.eps}}
	\caption{多图示例}
\end{figure}


\clearpage
\subsection{问题三分析}

题目要求制作软件的意思就是客户给定折叠桌高度、桌面边缘线的形状大小和桌脚边缘线的大致形状，将这些信息输入程序就得到客户想要的桌子。我们在求解最优设计加工参数时，自行给定桌面边缘线形状（椭圆、相交圆等），桌脚边缘线形状，折叠桌高度，应用第二问的非线性规划模型，用MATLAB软件绘制折叠桌截面图，得到自己设计的创意平板折叠桌。
